\renewcommand{\bibname}{Quellen und Bandzuordnung}
\begin{thebibliography}{99}
\addcontentsline{toc}{chapter}{Quellen und Bandzuordnung}
\footnotesize
\setlength{\itemsep}{2pt}
\setlength{\parskip}{0pt}
\interlinepenalty=10000

\bibitem{project}
Martin Kunze: \emph{Die Grundlagen der Mathematik}.
Band~3: die rein mengentheoretischen Kapitel ohne Auswahl.
Abgrenzung anhand des Überblicksbands, \path{VOLUMES.md} und
der Fachbandquellen; Format nach den Hauptresultate-Leitfäden
zu Band~28 und Band~37.

\bibitem{pila}
Jonathan Pila: \emph{B1.2 Set Theory. Lecture notes -- HT 2018}.
University of Oxford, 2018.
Abschnitt~2, Propositionen~2.6--2.8: Mengenbildung,
Durchschnitt und fehlende Allmenge; Abschnitte~3 und~7:
Potenzmenge, Ersetzung und Fundierung.
\url{https://www.maths.ox.ac.uk/system/files/attachments/SetTheoryHT18.pdf}.

\bibitem{koepke}
Peter Koepke: \emph{Set Theory 2006/07}.
Universität Bonn.
Abschnitte~2.4--2.7 zur Mengenbildung; daraus zusammengefaßt
das Mengenexistenzkriterium und das Kriterium für die Sammlung
aller Ergebnisse. Abschnitt~2.9 zur Fundierung.
\url{https://www.math.uni-bonn.de/ag/logik/teaching/2006WS/Script-2006WS-1.pdf}.

\bibitem{hammack}
Richard Hammack: \emph{Book of Proof}, dritte Auflage,
Fassung~3.4, 2018.
Abschnitte~1.3--1.6: Mengenoperationen;
Abschnitt~8.2, Beispiele~8.8--8.9: Potenzmengen-Inklusionen.
Die Lösungskriterien und die erweiterte Potenzmengendarstellung
sind hieraus und aus den Mengenbildungsaxiomen zusammengefaßt.
\url{https://richardhammack.github.io/BookOfProof/Main.pdf}.

\bibitem{paulson}
Lawrence C. Paulson und Mitwirkende:
\emph{Isabelle/ZF, Theory equalities}.
Abschnitte \emph{Set Difference} und
\emph{Big Union and Intersection}: allgemeine Mengengleichungen
und Inklusionen; insbesondere \texttt{Union\_Int\_subset}.
\url{https://www.cl.cam.ac.uk/research/hvg/Isabelle/dist/library/FOL/ZF/equalities.html}.

\bibitem{rutgers}
Rutgers University, MATH~215:
\emph{Set Operations Identities}, einseitige Formelsammlung.
Grundgesetze der Mengenoperationen.
\url{https://sites.math.rutgers.edu/~tb822/Sets%20Identities.pdf}.

\bibitem{cameron}
Peter J. Cameron: \emph{Notes on Algebraic Structures}, 2005/06.
Abschnitt~2.1, Punkt~6 und Proposition~2.3:
Mengengleichungen für symmetrische Differenz und Schnitt.
\url{https://webspace.maths.qmul.ac.uk/p.j.cameron/notes/algstr.pdf}.

\bibitem{miquel}
Alexandre Miquel: \emph{Preliminaries}, Kapitel~1 des
Forcing-Skripts, Universidad de la República.
Druckseiten~28--29, Aufgaben~1.9 und~1.13:
Mengeninduktion und Sammlungsschema in ZF.
\url{https://www.fing.edu.uy/~amiquel/forcing/forcing-chap1.pdf}.

\bibitem{holland}
James Holland: \emph{Set Theory Set Up}, Fassung 2023.
Abschnitt~9C, Theorem~9C\,2, Druckseite~99:
Sammlung mit der Zusatzbedingung für alle gesammelten Mengen.
\url{https://sites.math.rutgers.edu/~jch258/assets/Set%20Theory%20Set%20Up.pdf}.

\bibitem{kearnes}
Keith Kearnes: \emph{MATH 4730/5730: Set Theory},
University of Colorado Boulder, Vorlesungsübersicht vom
30.~August 2019: Durchschnitt einer beliebigen nichtleeren
definierbaren Gesamtheit von Mengen.
\url{https://math.colorado.edu/~kearnes/Notes/setsl.html}.

\end{thebibliography}
